\chapter{Sleep Medicine and Sleep Disorders}
\label{ch:sleep}
Sleep is a reversible behavioral state of perceptual disengagement from the environment, essential for physical restoration, cognitive processing, memory consolidation, metabolic homeostasis, and immune function. Sleep medicine encompasses the full spectrum of disorders affecting sleep architecture, circadian timing, breathing during sleep, excessive daytime sleepiness, abnormal movements, and undesirable behaviors during sleep. Nearly one-third of the adult population reports sleep complaints, with insomnia and obstructive sleep apnea representing the most prevalent disorders. Assessment of sleep disorders begins with a comprehensive sleep history covering bedtime routines, sleep latency, night awakenings, snoring, witnessed apneas, leg sensations, morning symptoms, daytime sleepiness (Epworth Sleepiness Scale), naps, caffeine and alcohol use, medications, and psychosocial stressors. Objective testing includes actigraphy, polysomnography (PSG), multiple sleep latency test (MSLT), and maintenance of wakefulness test (MWT).

%-------------------------------------------------------------
\section{Central Disorders of Hypersomnolence}
%-------------------------------------------------------------

%-------------------------------------------------------------

%-------------------------------------------------------------

%-------------------------------------------------------------

%-------------------------------------------------------------

%-------------------------------------------------------------

%-------------------------------------------------------------

%-------------------------------------------------------------

\label{sec:Central Disorders of Hypersomnolence}
%-------------------------------------------------------------%-------------------------------------------------------------

The central disorders of hypersomnolence are characterized by primary dysfunction of the central nervous system mechanisms regulating sleep and wakefulness, resulting in excessive daytime sleepiness (EDS) that is not attributable to other sleep disorders, medical conditions, or medication effects. Evaluation requires objective testing with polysomnography and the multiple sleep latency test (MSLT) to distinguish among narcolepsy type 1, narcolepsy type 2, idiopathic hypersomnia, and Kleine-Levin syndrome. MSLT is performed following an overnight polysomnogram of at least 6 hours, with five nap opportunities at 2-hour intervals measuring sleep latency and REM onset.

A mean sleep latency of $\le$8 minutes is considered pathologic. The presence of two or more sleep-onset REM periods (SOREMs) suggests narcolepsy. CSF hypocretin-1 measurement can confirm narcolepsy type 1 when clinical features are equivocal. HLA-DQB1*06:02 testing has limited diagnostic value due to high prevalence in the general population (12--25\%).

\subsection{Idiopathic Hypersomnia}
\label{sec:Idiopathic Hypersomnia}
Idiopathic hypersomnia is a disorder of excessive daytime sleepiness despite normal or prolonged nighttime sleep, with prolonged sleep inertia (severe difficulty awakening, confusion, grogginess, sometimes referred to as sleep drunkenness) and non-restorative naps. The prevalence is approximately 5--10 per 100,000. The condition results from unknown pathophysiology, with hypotheses including altered GABA-A receptor function, increased endogenous sleep-promoting substances, or dysfunction of wake-promoting circuits. Clinical features include excessive daytime sleepiness, prolonged sleep inertia, and non-restorative naps. Diagnosis requires MSLT showing mean sleep latency $\le$8 minutes with fewer than two SOREMs, or total 24-hour sleep time of >660 minutes on polysomnography or actigraphy, with exclusion of narcolepsy, sleep-disordered breathing, insufficient sleep syndrome, and other causes. Management includes wake-promoting agents (modafinil, armodafinil, pitolisant, solriamfetol), sodium oxybate, and clarithromycin. Prognosis is chronic; response to treatment is often less robust than in narcolepsy.

\subsection{Kleine-Levin Syndrome}
\label{sec:Kleine-Levin Syndrome}
Kleine-Levin syndrome (recurrent hypersomnia) is a rare disorder characterized by recurrent episodes of severe hypersomnia (16--20 hours of sleep per day), lasting days to weeks, with onset typically during adolescence. Prevalence is estimated at 1--2 per million, with a male predominance (2:1). The condition results from unknown pathophysiology, with proposed mechanisms including autoimmune or parainfectious inflammation in the hypothalamus. Clinical features include recurrent episodes of severe hypersomnia accompanied by hyperphagia, hypersexuality, mood changes, cognitive impairment, and sometimes psychotic features; between episodes, patients have normal sleep, cognition, and function. Diagnosis is clinical using ICSD-3 criteria. Management during episodes includes stimulants (modafinil, methylphenidate) for sleepiness, lithium for episode prevention, and mood stabilizers (valproate, lamotrigine). Prognosis is favorable; the disease typically resolves spontaneously over 8--14 years.

\subsection{Narcolepsy Type 1}
\label{sec:Narcolepsy Type 1}
Narcolepsy type 1 (narcolepsy with cataplexy) is a chronic neurological disorder characterized by excessive daytime sleepiness, cataplexy (sudden, transient episodes of bilateral skeletal muscle weakness triggered by strong emotions), sleep paralysis, hypnagogic hallucinations, and fragmented nighttime sleep. The prevalence is approximately 25--50 per 100,000 individuals, with onset peaking in adolescence and young adulthood. The condition results from autoimmune-mediated destruction of hypocretin (orexin)-producing neurons in the lateral hypothalamus, leading to severely decreased or absent CSF hypocretin-1 levels; genetic susceptibility is strongly linked to HLA-DQB1*06:02 (present in >90\% of cases). Clinical features include excessive daytime sleepiness, cataplexy (brief episodes of bilateral weakness provoked by laughter or surprise), sleep paralysis, and hypnagogic hallucinations. Diagnosis requires polysomnography followed by MSLT demonstrating mean sleep latency $\le$8 minutes and two or more SOREM naps, plus either cataplexy or CSF hypocretin-1 deficiency. Management includes scheduled naps, wake-promoting agents (modafinil, armodafinil, pitolisant, solriamfetol), anti-cataplexy agents (sodium oxybate, antidepressants), and sodium oxybate for both EDS and cataplexy. Prognosis is chronic but symptoms can be controlled with appropriate therapy.

\subsection{Narcolepsy Type 2}
\label{sec:Narcolepsy Type 2}
Narcolepsy type 2 is a chronic neurological disorder characterized by excessive daytime sleepiness and objective findings on MSLT (mean sleep latency $\le$8 minutes with two or more SOREMs) but without cataplexy and with normal CSF hypocretin-1 levels. It represents 15--25\% of narcolepsy cases. Clinical features are generally less severe than type 1; some patients may eventually develop cataplexy and transition to type 1. Diagnosis requires polysomnography followed by MSLT demonstrating the characteristic findings without cataplexy. Management focuses on wake-promoting agents (modafinil, armodafinil, pitolisant, solriamfetol) and behavioral sleep hygiene. Prognosis is chronic but manageable.

%-------------------------------------------------------------
\section{Circadian Rhythm Sleep-Wake Disorders}
%-------------------------------------------------------------

%-------------------------------------------------------------

%-------------------------------------------------------------

%-------------------------------------------------------------

%-------------------------------------------------------------

%-------------------------------------------------------------

%-------------------------------------------------------------

%-------------------------------------------------------------

\label{sec:Circadian Rhythm Sleep-Wake Disorders}
%-------------------------------------------------------------%-------------------------------------------------------------

Circadian rhythm sleep-wake disorders (CRSWDs) result from a persistent or recurrent pattern of sleep disruption caused by alteration of the endogenous circadian timing system or misalignment between the internal circadian rhythm and the external environment. Patients present with insomnia, excessive daytime sleepiness, or both, with significant distress or impairment. Diagnosis requires at least 7 days (preferably 14 days) of actigraphy and sleep diaries to characterize the timing and pattern of sleep. Dim-light melatonin onset (DLMO) is the gold-standard biomarker for circadian phase assessment.

\subsection{Advanced Sleep-Wake Phase Disorder}
\label{sec:Advanced Sleep-Wake Phase Disorder}
Advanced sleep-wake phase disorder (ASWPD) is a circadian rhythm sleep-wake disorder characterized by sleep-wake timing that is several hours earlier than desired, with sleep onset between 6:00--9:00 PM and spontaneous awakening between 2:00--5:00 AM. It is less common than DSWPD, affecting approximately 1\% of the population, predominantly older adults (though familial cases with PER2 or PER3 gene mutations exist). The condition results from a short intrinsic circadian period (<24 hours) or an exaggerated phase-advance response to morning light. Clinical features include early evening sleepiness and early morning awakening. Diagnosis is by sleep diaries and actigraphy. Management involves bright light therapy in the evening to phase-delay the circadian clock and low-dose melatonin upon early morning awakening. Prognosis is favorable with treatment compliance.

\subsection{Delayed Sleep-Wake Phase Disorder}
\label{sec:Delayed Sleep-Wake Phase Disorder}
Delayed sleep-wake phase disorder (DSWPD) is a circadian rhythm sleep-wake disorder characterized by a stable delay of the sleep-wake phase relative to the desired or required schedule, with sleep onset typically after 1:00--3:00 AM and wake times in the late morning or early afternoon. Prevalence is 1--3\% in adults and 7--16\% in adolescents, making it the most common circadian rhythm disorder. The condition results from an abnormally long intrinsic circadian period, reduced phase-advance capacity, or a hypo-responsive phase-advance mechanism to morning light. Clinical features include severe difficulty falling asleep at conventional times but normal sleep once initiated when allowed to follow the endogenous schedule. Diagnosis is by clinical history, sleep diaries, actigraphy, and dim-light melatonin onset measurement. Management includes timed bright light therapy in the morning, low-dose melatonin (0.5--3 mg) administered 5--6 hours before target bedtime, and consistent sleep-wake scheduling; chronotherapy may be considered in refractory cases. Prognosis is favorable with treatment compliance.

\subsection{Irregular Sleep-Wake Rhythm Disorder}
\label{sec:Irregular Sleep-Wake Rhythm Disorder}
Irregular sleep-wake rhythm disorder (ISWRD) is a circadian rhythm sleep-wake disorder presenting as a fragmented, disorganized sleep-wake pattern with at least three distinct sleep episodes within a 24-hour period, without a clear circadian rhythm. It is most common in patients with neurodegenerative disorders (Alzheimer disease, dementia), traumatic brain injury, intellectual disability, and in institutionalized elderly populations. The condition results from degeneration of the SCN or reduced photic input. Clinical features include multiple sleep episodes throughout the day and night with no dominant sleep period. Diagnosis requires actigraphy or sleep diaries over at least 7 days. Management focuses on increasing daytime bright light exposure, structured daytime activities and social schedules, evening melatonin, and avoidance of nighttime light; combination therapy (bright light plus melatonin plus structured activity) is most effective. Prognosis is variable and depends on the underlying condition.

\subsection{Jet Lag Disorder}
\label{sec:Jet Lag Disorder}
Jet lag disorder is a temporary circadian misalignment caused by rapid travel across multiple time zones. Symptoms include difficulty initiating or maintaining sleep at the destination bedtime, daytime fatigue, impaired cognitive performance, gastrointestinal disturbances, and general malaise. Severity correlates with the number of time zones crossed (typically three or more) and direction of travel (eastward travel is more difficult because it requires phase-advancing, which is harder than phase-delaying). The condition results from the circadian system adapting at a rate of approximately one time zone per day. Clinical features include sleep disturbance and daytime impairment following rapid transmeridian travel. Diagnosis is clinical. Management includes controlled light exposure (seek morning light when traveling east, evening light when traveling west), melatonin (0.5--5 mg at destination bedtime), and short-term use of hypnotics for sleep initiation. Prognosis is favorable; symptoms resolve as the circadian system adapts.

\subsection{Non-24-Hour Sleep-Wake Rhythm Disorder}
\label{sec:Non-24-Hour Sleep-Wake Rhythm Disorder}
Non-24-hour sleep-wake rhythm disorder (N24) is a circadian rhythm sleep-wake disorder characterized by a sleep-wake cycle that is desynchronized from the 24-hour day, typically running at a period longer than 24 hours (free-running), with each day sleep onset and wake time drifting progressively later by 30--60 minutes. It occurs in 50--80\% of blind individuals (due to lack of photic entrainment of the SCN) and rarely in sighted individuals. The condition results from absence of light perception to entrain the circadian clock. Clinical features include a progressive daily delay in sleep-wake timing cycling through the 24-hour day. Diagnosis requires at least 14 days of actigraphy or sleep diaries demonstrating a circadian period other than 24 hours. Management in blind individuals includes melatonin (0.5--10 mg at desired bedtime) and tasimelteon (approved for N24 in blind people); in sighted individuals, treatment includes morning light therapy and evening melatonin. Prognosis is chronic but manageable.

\subsection{Shift Work Sleep Disorder}
\label{sec:Shift Work Sleep Disorder}
Shift work sleep disorder (SWSD) is a circadian rhythm sleep-wake disorder affecting individuals whose work schedules overlap with the typical sleep period (night shifts, early morning shifts, rotating shifts), causing circadian misalignment with resultant insomnia, excessive sleepiness, or both. Approximately 20--30\% of shift workers meet criteria. The condition results from a misalignment between the endogenous circadian timing system and the work-rest schedule. Clinical features include difficulty initiating or maintaining sleep during the available sleep period and excessive sleepiness during work hours. Diagnosis is clinical with sleep diaries. Management strategies include strategic napping before or during shifts, bright light exposure during work, controlled light avoidance during the commute home, fixed sleep scheduling in a dark/cool room, melatonin at the desired bedtime, and wake-promoting agents (modafinil, armodafinil, or caffeine) for night-shift alertness. Prognosis is favorable with appropriate management; chronic misalignment increases risks of cardiovascular disease, metabolic syndrome, and depression.

%-------------------------------------------------------------
\section{Insomnia}
%-------------------------------------------------------------

%-------------------------------------------------------------

%-------------------------------------------------------------

%-------------------------------------------------------------

Insomnia is a prevalent sleep disorder characterized by persistent difficulty initiating or maintaining sleep, or experiencing non-restorative sleep despite adequate opportunity. Neurobiological hyperarousal involving heightened sympathetic tone and hypothalamic-pituitary-adrenal axis activation disrupts normal sleep-wake regulation. Cognitive behavioral therapy for insomnia (CBT-I) represents first-line therapy, supported by short-term pharmacotherapy when appropriate.

%-------------------------------------------------------------

%-------------------------------------------------------------

%-------------------------------------------------------------

%-------------------------------------------------------------

\label{sec:Insomnia}
%-------------------------------------------------------------%-------------------------------------------------------------

\subsection{Acute/Short-Term Insomnia}
\label{sec:Acute Short-Term Insomnia}
Acute insomnia is a sleep disturbance lasting less than three months, typically triggered by a specific stressor (acute medical illness, hospitalization, surgery, bereavement, work stress, jet lag). It is extremely common, affecting 30--50\% of adults annually. The condition results from transient hyperarousal and maladaptive behavioral responses (extended time in bed, napping, irregular sleep schedules). Clinical features include difficulty initiating or maintaining sleep that is temporally linked to an identifiable stressor. Diagnosis is clinical. Management involves sleep hygiene, stimulus control, and short-term pharmacotherapy limited to 1--2 weeks to prevent progression to chronic insomnia. Prognosis is favorable; most cases resolve spontaneously when the stressor resolves.

\subsection{Chronic Insomnia Disorder}
\label{sec:Chronic Insomnia Disorder}
Chronic insomnia disorder is a sleep disorder defined by persistent difficulty initiating sleep, maintaining sleep, or early morning awakening with inability to return to sleep, despite adequate opportunity for sleep, occurring at least three nights per week for at least three months, with significant daytime impairment (fatigue, mood disturbance, cognitive dysfunction, or reduced performance). The prevalence is 6--10\% in the general population, with females affected 1.5--2 times more frequently than males. Risk factors include advanced age, female sex, shift work, medical and psychiatric comorbidities, and genetic predisposition. The condition results from hyperarousal of the hypothalamic-pituitary-adrenal axis (elevated cortisol, ACTH), increased high-frequency EEG activity during sleep, activation of the sympathetic nervous system, and maladaptive conditioning (negative associations with the sleep environment). Clinical features include difficulty falling asleep, frequent awakenings, early morning awakening, and daytime impairment. Diagnosis is clinical based on ICSD-3 and DSM-5 criteria, supplemented by sleep diaries and actigraphy; polysomnography is not routinely indicated but may evaluate comorbid sleep disorders. Management involves cognitive behavioral therapy for insomnia (CBT-I) as first-line therapy---stimulus control, sleep restriction, cognitive restructuring, and sleep hygiene education---with pharmacotherapy (melatonin, ramelteon, orexin receptor antagonists, benzodiazepine receptor agonists, trazodone, doxepin) as adjunctive treatment. Prognosis is favorable with CBT-I; long-term outcomes favor CBT-I over medication alone.

\subsection{Idiopathic Insomnia}
\label{sec:Idiopathic Insomnia}
Idiopathic insomnia is a lifelong form of insomnia with onset in infancy or early childhood, without identifiable triggers, and in the absence of medical, psychiatric, or environmental causes. It is rare, representing less than 5\% of insomnia cases. The condition results from a hypothesized primary abnormality in the sleep-wake regulatory system, such as dysfunction of the VLPO or an imbalance between sleep-promoting and arousal-promoting systems. Clinical features include objectively reduced total sleep time on polysomnography and persistent hyperarousal. Diagnosis is clinical, with polysomnography used to rule out other sleep disorders. Management involves CBT-I as first-line therapy, though efficacy may be limited; pharmacotherapy (hypnotics, sedating antidepressants) is often required long-term. Prognosis is persistence across the lifespan.

\subsection{Insomnia in the Context of Medical/Psychiatric Disorders}
\label{sec:Insomnia in the Context of Medical Psychiatric Disorders}
Insomnia comorbid with medical or psychiatric conditions is a sleep disorder that represents the most common presentation of insomnia in clinical practice. Psychiatric disorders---particularly major depressive disorder, generalized anxiety disorder, panic disorder, and post-traumatic stress disorder---are the strongest risk factors, with bidirectional relationships. The condition results from the underlying medical or psychiatric disorder and its effects on sleep regulation. Clinical features include difficulty initiating or maintaining sleep in the presence of a comorbid condition. Diagnosis is clinical, with identification of the underlying medical or psychiatric disorder. Management involves treatment of the underlying condition, CBT-I adapted for the specific comorbidity, and cautious use of hypnotics mindful of drug-disease and drug-drug interactions; depression with comorbid insomnia may benefit from sedating antidepressants (trazodone, mirtazapine, amitriptyline). Prognosis depends on the course of the underlying condition.

\subsection{Psychophysiologic Insomnia}
\label{sec:Psychophysiologic Insomnia}
Psychophysiologic insomnia is a subtype of chronic insomnia characterized by a learned, conditioned arousal to sleep-related stimuli, such that the bed and bedroom environment become associated with frustration, worry, and heightened alertness rather than sleep. Patients often sleep better away from home (paradoxical improvement in novel environments). The condition results from physiologic hyperarousal with elevated metabolic rate, heart rate, and cortisol at sleep onset, combined with cognitive patterns including excessive preoccupation with sleep, catastrophic thinking about sleep loss, and performance anxiety. Clinical features include difficulty sleeping in one's own bed but improved sleep in novel environments. Diagnosis is clinical. Management involves stimulus control therapy (reassociating the bed with sleep), sleep restriction therapy (consolidating sleep window to improve sleep efficiency), and cognitive restructuring to reduce maladaptive beliefs about sleep. Prognosis is favorable with appropriate behavioral therapy.

%-------------------------------------------------------------
\section{Parasomnias}
%-------------------------------------------------------------

%-------------------------------------------------------------

%-------------------------------------------------------------

%-------------------------------------------------------------

%-------------------------------------------------------------

%-------------------------------------------------------------

%-------------------------------------------------------------

%-------------------------------------------------------------

\label{sec:Parasomnias}
%-------------------------------------------------------------%-------------------------------------------------------------

Parasomnias are undesirable physical events or experiences that occur during entry into sleep, within sleep, or during arousal from sleep. They are categorized by the sleep stage of origin: NREM-related parasomnias (disorders of arousal), REM-related parasomnias, and other parasomnias. Diagnosis relies on detailed clinical history from both the patient and bed partner, supported by polysomnography with synchronized video recording (vPSG) when episodes are frequent, potentially injurious, atypical, or when differentiation from nocturnal seizures is required.

\subsection{NREM Parasomnias: Confusional Arousals}
\label{sec:Confusional Arousals}
Confusional arousals are episodes of mental confusion, disorientation, and inappropriate behavior (slow speech, fumbling, ineffective actions) upon arousal from NREM sleep, typically N3. Episodes are more common in children. The condition results from incomplete arousal from slow-wave sleep. Clinical features include confusion and disorientation upon awakening without the complex mobility of sleepwalking or the intense autonomic activation of sleep terrors. Severe confusional arousal in adults is associated with obstructive sleep apnea, shift work, and psychiatric illness. Diagnosis is clinical. Management addresses underlying sleep-disordered breathing and triggers; clonazepam may be used in severe cases. Prognosis is favorable.

\subsection{NREM Parasomnias: Sleep Terrors}
\label{sec:Sleep Terrors}
Sleep terrors (pavor nocturnus) are dramatic arousals from NREM slow-wave sleep, characterized by sudden sitting up, piercing screams, intense autonomic activation (tachycardia, tachypnea, mydriasis, sweating), and behavioral expressions of intense fear. Peak prevalence is in children ages 1.5--7 years (1--6\%). The condition results from incomplete arousal from deep NREM sleep. Clinical features include episodes lasting 1--10 minutes with the individual inconsolable, confused, and potentially thrashing violently, followed by return to sleep with no recall of the event. Diagnosis is clinical. Management includes reassurance, parental education, safety measures, and clonazepam for severe cases. Prognosis is favorable; most children outgrow the condition.

\subsection{NREM Parasomnias: Sleep-Related Eating Disorder}
\label{sec:Sleep-Related Eating Disorder}
Sleep-related eating disorder (SRED) is a disorder involving recurrent episodes of involuntary eating and drinking during partial arousals from NREM sleep, with impaired awareness and variable amnesia for the events. Prevalence is 1--3\% in the general population, higher in women. The condition results from abnormal arousal from NREM sleep and is strongly associated with other NREM parasomnias and hypnotic use (particularly zolpidem). Clinical features include consumption of unusual or dangerous items during the night with subsequent weight gain, dental injury, and potential ingestion of toxic substances. Diagnosis is clinical. Management includes discontinuation of triggering hypnotics, treatment of comorbid sleep apnea or restless legs syndrome, topiramate (50--200 mg nightly), and safety measures. Prognosis is favorable with appropriate treatment.

\subsection{NREM Parasomnias: Sleepwalking}
\label{sec:Sleepwalking}
Sleepwalking (somnambulism) is a disorder of arousal from NREM sleep, typically arising from N3 (slow-wave) sleep during the first third of the night. Prevalence is 10--30\% in children and 2--4\% in adults. The condition results from incomplete arousal from slow-wave sleep, with genetic factors playing a strong role. Clinical features include complex, semi-purposeful behaviors ranging from sitting up to walking and performing routine activities without conscious awareness, with a blank stare and amnesia for the event. Triggers include sleep deprivation, fever, alcohol, stress, and sedative-hypnotic use. Diagnosis is clinical; polysomnography is indicated for atypical, frequent, or potentially injurious episodes. Management includes safety precautions, avoidance of triggers, scheduled awakening before typical episode time, and pharmacotherapy (clonazepam, diazepam) for frequent or injurious episodes. Prognosis is favorable; most children outgrow the condition.

\subsection{Other Parasomnias: Exploding Head Syndrome}
\label{sec:Exploding Head Syndrome}
Exploding head syndrome is a parasomnia characterized by episodes of loud imagined noises (explosion, gunshot, crashing cymbals) or a sensation of violent explosion inside the head, occurring at the transition from wakefulness to sleep (hypnagogic) or from sleep to wakefulness. It is relatively common (10--15\% of young adults). The condition results from a transient delay in the normal inhibition of the auditory system during sleep onset. Clinical features include frightening but painless auditory sensations lasting seconds. Diagnosis is clinical. Management involves reassurance; clomipramine and calcium channel blockers have been used in severe cases. Prognosis is benign.

\subsection{Other Parasomnias: Sleep Enuresis}
\label{sec:Sleep Enuresis}
Sleep enuresis (bedwetting) is involuntary voiding of urine during sleep, occurring at least twice per week for at least three months, in children over age 5. Prevalence decreases with age: 15\% at age 5, 5\% at age 10, and 1--2\% by age 15, with a male predominance. The condition results from a mismatch between nocturnal bladder capacity and overnight urine production, impaired arousal from sleep in response to bladder fullness, and in some cases, nocturnal polyuria due to insufficient antidiuretic hormone secretion. Primary enuresis (never dry for an extended period) is more common than secondary enuresis (relapse after $\ge$6 months of dryness). Diagnosis includes urinalysis to exclude diabetes mellitus, diabetes insipidus, and urinary tract infection. Management includes behavioral interventions (voiding before bed, limiting evening fluids, reward systems, bedwetting alarms), desmopressin for nocturnal polyuria, and anticholinergics for bladder overactivity; combination therapy may benefit refractory cases. Prognosis is favorable; most children outgrow the condition.

\subsection{Other Parasomnias: Sleep-Related Hallucinations}
\label{sec:Sleep-Related Hallucinations}
Sleep-related hallucinations are vivid, dream-like perceptual experiences occurring at sleep onset (hypnagogic) or upon awakening (hypnopompic), involving visual, auditory, tactile, or kinetic sensations. They are common in the general population (25--40\% of young adults) and typically benign. The condition results from intrusion of REM sleep imagery into wakefulness; pathological sleep-related hallucinations are hallmark features of narcolepsy type 1 and can occur in Parkinson disease. Clinical features include multimodal hallucinations with intact reality testing. Diagnosis requires exclusion of narcolepsy and distinction from psychotic disorders. Management involves reassurance for benign cases and treatment of underlying narcolepsy or neurodegenerative disease when present. Prognosis is benign in isolation.

\subsection{REM Parasomnias: Nightmare Disorder}
\label{sec:Nightmare Disorder}
Nightmare disorder is a parasomnia involving recurrent, vividly realistic, dysphoric dreams that awaken the individual from REM sleep, typically during the second half of the night, with rapid return to full alertness and clear recall of the frightening content. Prevalence is 2--8\% in adults and 10--50\% in children. The condition results from abnormal processing of emotional memories during REM sleep and is associated with PTSD, anxiety, borderline personality disorder, and certain medications (beta-blockers, SSRIs). Clinical features include recurrent nightmares causing sleep avoidance and daytime fatigue. Diagnosis is clinical. Management includes imagery rehearsal therapy (IRT), prazosin (an alpha-1 adrenergic antagonist) for PTSD-related nightmares, and treatment of comorbid psychiatric disorders. Prognosis is favorable with treatment.

\subsection{REM Parasomnias: Recurrent Isolated Sleep Paralysis}
\label{sec:Recurrent Isolated Sleep Paralysis}
Recurrent isolated sleep paralysis is a parasomnia consisting of episodes of inability to perform voluntary movements at sleep onset (hypnagogic) or upon awakening (hypnopompic), with preserved consciousness and awareness of the environment. Prevalence is 8--50\% depending on population, higher in those with anxiety, panic disorder, and sleep deprivation. The condition results from a dissociation between REM sleep atonia and conscious awareness. Clinical features include episodes lasting seconds to minutes, often accompanied by terrifying hallucinations (intruder, pressure on the chest), and the ability to terminate episodes by external stimulation or strong voluntary effort to move a small muscle group. Diagnosis is clinical. Management includes reassurance and sleep hygiene. Prognosis is benign.

\subsection{REM Parasomnias: REM Sleep Behavior Disorder}
\label{sec:REM Sleep Behavior Disorder}
REM sleep behavior disorder (RBD) is a parasomnia characterized by loss of the normal muscle atonia of REM sleep, allowing patients to act out their dreams, often with violent or injurious behaviors. Prevalence is 1--2\% in the general population, with a strong male predominance (>80\% of cases), typically occurring after age 50. The condition results from dysfunction of the brainstem REM-atonia circuits (sublaterodorsal nucleus in the pons). Critically, idiopathic RBD is a highly specific prodromal marker for synucleinopathies: 80--90\% of patients will eventually develop Parkinson disease, dementia with Lewy bodies, or multiple system atrophy, with a mean latency of 10--15 years. Clinical features include dream-enactment behaviors (punching, kicking, jumping out of bed) corresponding to aggressive dream content. Diagnosis is clinical with polysomnography demonstrating elevated chin EMG tone during REM. Management includes safety measures, clonazepam (first-line, 0.5--2 mg nightly), and melatonin (3--15 mg) as an effective alternative. Prognosis is chronic; the condition is manageable but carries a high risk of eventual neurodegenerative disease.

%-------------------------------------------------------------
\section{Physiology of Sleep}
%-------------------------------------------------------------

%-------------------------------------------------------------

%-------------------------------------------------------------

%-------------------------------------------------------------

%-------------------------------------------------------------

%-------------------------------------------------------------

%-------------------------------------------------------------

%-------------------------------------------------------------

\label{sec:Physiology of Sleep}
%-------------------------------------------------------------%-------------------------------------------------------------

\subsection{Circadian Rhythm Regulation}
\label{sec:Circadian Rhythm Regulation}
Circadian rhythms are endogenous near-24-hour biological cycles regulated by the suprachiasmatic nucleus (SCN) of the anterior hypothalamus. The SCN receives photic input from intrinsically photosensitive retinal ganglion cells (ipRGCs containing melanopsin) via the retinohypothalamic tract and coordinates peripheral clocks throughout the body via neuroendocrine (melatonin, cortisol) and autonomic outputs. Melatonin, synthesized and released by the pineal gland during darkness, is the principal hormonal signal for sleep onset and phase shifting. Cortisol exhibits a reciprocal pattern, with nadir around midnight and peak upon awakening (cortisol awakening response). The circadian system interacts homeostatically with the sleep-wake drive (Process C and Process S): adenosine accumulation during wakefulness drives homeostatic sleep pressure (Process S), while circadian alerting signals (Process C) counterbalance this pressure to maintain consolidated wakefulness. The two-process model explains the timing and intensity of sleep under normal and sleep-deprived conditions. Circadian misalignment occurs when the endogenous clock is desynchronized from the external light-dark cycle or social schedule.

\subsection{Sleep Architecture Across the Lifespan}
\label{sec:Sleep Architecture Across the Lifespan}
Sleep architecture is the pattern of sleep stages that changes dramatically across the lifespan. Newborns spend approximately 50\% of sleep time in active (REM-like) sleep with an ultradian cycle of 50--60 minutes, and total sleep duration averages 16--18 hours daily. By age 5, REM proportion declines to adult levels (20--25\%), and N3 sleep reaches its peak in childhood. During adolescence, a physiologic phase delay emerges (delayed melatonin onset), coupled with reduced slow-wave sleep. Young adults exhibit well-defined sleep cycles with 20--25\% REM, 45--55\% N2, and 15--25\% N3. Aging is associated with progressive decline in slow-wave sleep (especially N3), increased sleep fragmentation, increased N1 and wake after sleep onset (WASO), decreased total sleep time, phase advancement of circadian rhythms, and reduced amplitude of the melatonin rhythm. The prevalence of sleep disorders such as insomnia, sleep-disordered breathing, and REM sleep behavior disorder increases with age.

\subsection{Sleep Stages (NREM N1--N3, REM)}
\label{sec:Sleep Stages}
Sleep is a reversible behavioral state divided into non-rapid eye movement (NREM) and rapid eye movement (REM) sleep, alternating in approximately 90-minute cycles throughout the night. NREM sleep comprises three stages. N1 (Stage 1) is the lightest stage, representing the transition from wakefulness, characterized by theta rhythm (4--7 Hz) on EEG, slow rolling eye movements, and easy arousal. N2 (Stage 2) is the predominant sleep stage, occupying approximately 45--55\% of total sleep time in adults, characterized by sleep spindles (12--14 Hz bursts) and K-complexes on EEG with decreased heart rate and core body temperature. N3 (Stage 3), or slow-wave sleep, is characterized by high-amplitude, low-frequency delta waves (0.5--4 Hz), a high arousal threshold, and the greatest anabolic and restorative functions, including growth hormone secretion, immune cytokine release, and synaptic downscaling. REM sleep is characterized by rapid, conjugate eye movements, skeletal muscle atonia (except diaphragm and extraocular muscles), desynchronized low-voltage mixed-frequency EEG, and vivid dreaming. Autonomic instability during REM produces fluctuating heart rate, blood pressure, and respiration. The transition from wakefulness to sleep involves activation of the ventrolateral preoptic nucleus (VLPO) in the hypothalamus, which inhibits arousal-promoting brainstem nuclei (locus coeruleus, raphe nuclei, tuberomammillary nucleus). The flip-flop switch model describes the mutual inhibition between the VLPO and the monoaminergic arousal systems, ensuring rapid and stable state transitions.

%-------------------------------------------------------------
\section{Sleep-Related Breathing Disorders}
%-------------------------------------------------------------

%-------------------------------------------------------------

%-------------------------------------------------------------

%-------------------------------------------------------------

%-------------------------------------------------------------

%-------------------------------------------------------------

%-------------------------------------------------------------

%-------------------------------------------------------------

\label{sec:Sleep-Related Breathing Disorders}
%-------------------------------------------------------------%-------------------------------------------------------------

Sleep-related breathing disorders encompass a spectrum of conditions characterized by abnormal respiration during sleep, including obstructive sleep apnea, central sleep apnea syndromes, sleep-related hypoventilation disorders, and sleep-related hypoxemia disorder. These conditions produce intermittent hypoxia, hypercapnia, sleep fragmentation, and autonomic dysregulation, contributing to significant cardiovascular, metabolic, and neurocognitive morbidity. Diagnosis requires polysomnography with detailed respiratory monitoring (oronasal thermistor, nasal pressure transducer, thoracic and abdominal respiratory inductance plethysmography, pulse oximetry, and capnography).

\subsection{Central Sleep Apnea}
\label{sec:Central Sleep Apnea}
Central sleep apnea (CSA) is a sleep-related breathing disorder characterized by recurrent episodes of absent or reduced respiratory effort during sleep due to instability of the central respiratory control system, without airway obstruction. CSA affects 5--10\% of patients referred for sleep studies and is frequently comorbid with heart failure and neurological conditions. The condition results from increased chemosensitivity (high loop gain), prolonged circulation time (Cheyne-Stokes respiration in heart failure), hypobaric hypoxia-induced hyperventilation and hypocapnia (high-altitude periodic breathing), opioid-induced respiratory depression, or brainstem lesions. Clinical features include fragmented sleep, nocturnal dyspnea, and excessive daytime sleepiness. Diagnosis requires polysomnography demonstrating five or more central apneas or hypopneas per hour with no evidence of obstruction. Management involves treatment of underlying conditions, supplemental oxygen, positive airway pressure (CPAP, bilevel, or adaptive servo-ventilation), acetazolamide, and gradual acclimatization for high-altitude periodic breathing. Prognosis depends on the underlying cause; optimization of cardiac function is essential in heart failure patients.

\subsection{Complex Sleep Apnea}
\label{sec:Complex Sleep Apnea}
Complex sleep apnea (treatment-emergent central sleep apnea) is a sleep-related breathing disorder defined by the emergence or persistence of CSA during PAP therapy in patients initially diagnosed with OSA. It occurs in 5--15\% of patients initiating CPAP therapy, particularly those with cardiovascular disease, opioid use, or high loop gain physiology. The condition results from PAP-induced elimination of the obstructive component, unmasking a latent central ventilatory instability. Clinical features include persistent respiratory events despite adequate PAP therapy for OSA. Diagnosis requires polysomnography on PAP demonstrating central events. Management involves continued PAP adherence, as many cases resolve spontaneously within 3--6 months; persistent cases may require switching to adaptive servo-ventilation or bilevel PAP with a backup rate. Prognosis is generally favorable.

\subsection{Obesity Hypoventilation Syndrome}
\label{sec:Obesity Hypoventilation Syndrome}
Obesity hypoventilation syndrome (OHS, Pickwickian syndrome) is a sleep-related breathing disorder defined by the triad of obesity (BMI $\ge$30 kg/m$^2$), awake daytime hypercapnia (PaCO$_2$ >45 mmHg) without other causes, and sleep-disordered breathing (90\% have OSA). Prevalence among obese individuals is 10--20\%, with higher rates in severe obesity. The condition results from mechanical loading of the respiratory system (reduced chest wall compliance, diaphragmatic restriction), blunted central respiratory drive (decreased chemosensitivity to hypercapnia and hypoxia), and leptin resistance leading to impaired ventilatory stimulation. Clinical features include obesity, hypersomnolence, dyspnea on exertion, morning headache, cor pulmonale (right heart failure), and pulmonary hypertension. Diagnosis requires arterial blood gas showing awake hypercapnia, polysomnography, and exclusion of other causes of hypoventilation (neuromuscular disease, COPD, hypothyroidism). Management involves positive airway therapy (CPAP or bilevel PAP), weight loss (medical or bariatric surgery), and supplemental oxygen if needed. Prognosis is guarded; untreated OHS is associated with significantly increased cardiovascular morbidity and mortality.

\subsection{Obstructive Sleep Apnea}
\label{sec:Obstructive Sleep Apnea}
Obstructive sleep apnea (OSA) is a sleep-related breathing disorder characterized by recurrent episodes of partial or complete pharyngeal collapse during sleep, resulting in apnea (breathing cessation $\ge$10 seconds) or hypopnea (reduced airflow with desaturation or arousal), causing intermittent hypoxia, hypercapnia, intrathoracic pressure swings, and sleep fragmentation. Prevalence is 15--30\% of men and 10--20\% of women in the general population, with marked increases in obesity. The condition results from anatomic narrowing of the pharyngeal airway (obesity, craniofacial abnormalities, macroglossia, tonsillar hypertrophy) combined with reduced dilator muscle activity of the genioglossus during sleep and a low arousal threshold. Clinical features include loud snoring, witnessed apneas, gasping or choking awakenings, nocturia, morning headache, excessive daytime sleepiness, fatigue, cognitive impairment, and irritability. Diagnosis is by attended in-laboratory polysomnography (gold standard) or home sleep apnea testing (HSAT) in uncomplicated adults, with severity classified by the apnea-hypopnea index (AHI): mild (AHI 5--15/hr), moderate (AHI 15--30/hr), and severe (AHI >30/hr). Management involves positive airway pressure (PAP) therapy (continuous PAP, auto-CPAP, bilevel PAP) as first-line treatment, mandibular advancement devices for mild-to-moderate OSA, hypoglossal nerve stimulation therapy, positional therapy, weight loss and lifestyle modification, and upper airway surgery for patients with surgically correctable anatomy or PAP intolerance. Prognosis is favorable with treatment; untreated OSA is associated with hypertension, atrial fibrillation, heart failure, coronary artery disease, stroke, and type 2 diabetes.

%-------------------------------------------------------------
\section{Sleep-Related Movement Disorders}
%-------------------------------------------------------------

%-------------------------------------------------------------

%-------------------------------------------------------------

%-------------------------------------------------------------

%-------------------------------------------------------------

%-------------------------------------------------------------

%-------------------------------------------------------------

%-------------------------------------------------------------

\label{sec:Sleep-Related Movement Disorders}
%-------------------------------------------------------------%-------------------------------------------------------------

Sleep-related movement disorders are characterized by relatively simple, usually stereotyped movements that disturb sleep or its onset. These conditions are distinct from parasomnias in that the movements are typically simple and repetitive rather than complex behaviors. Diagnosis relies on clinical history, and polysomnography with bilateral anterior tibialis EMG (and additional muscle groups as indicated) provides objective confirmation.

\subsection{Periodic Limb Movement Disorder}
\label{sec:Periodic Limb Movement Disorder}
Periodic limb movement disorder (PLMD) is a sleep-related movement disorder characterized by recurrent, stereotypical, involuntary movements of the lower extremities during sleep, most commonly involving dorsiflexion of the great toe, fanning of the toes, and flexion of the ankle, knee, or hip, occurring at intervals of 5--90 seconds. Prevalence increases with age, affecting 10--30\% of adults over 60, but is much lower as a primary diagnosis. The condition results from unknown pathophysiology, likely related to dopaminergic dysfunction. Clinical features include sleep fragmentation from movements associated with microarousals, leading to daytime impairment. Diagnosis requires polysomnography demonstrating a periodic limb movement index (PLMI) >15 per hour in adults (>5 in children) with clinically significant sleep disturbance, in the absence of RLS or other explanatory disorders. Management is indicated only when PLMS are associated with sleep disturbance and includes alpha-2-delta ligands (gabapentin, pregabalin) preferred over dopamine agonists due to augmentation risk. Prognosis is chronic but manageable.

\subsection{Propriospinal Myoclonus at Sleep Onset}
\label{sec:Propriospinal Myoclonus at Sleep Onset}
Propriospinal myoclonus at sleep onset is a rare sleep-related movement disorder characterized by sudden, brief, involuntary jerks of the axial muscles occurring during the transition from wakefulness to sleep. The condition results from myoclonus originating in the spinal cord and propagating rostrally and caudally via propriospinal pathways. Clinical features include jerks triggered by voluntary attempts to relax or fall asleep, severely delaying sleep onset and resulting in sleep-onset insomnia. Diagnosis requires polygraphic EMG recording across multiple spinal segments and polysomnography, with differentiation from hypnic jerks (physiologic sleep starts) and epileptic myoclonus. Management is challenging; clonazepam, levetiracetam, and other anticonvulsants have been used with variable success. Prognosis is variable; some patients experience spontaneous improvement, while others require long-term pharmacotherapy.

\subsection{Restless Legs Syndrome}
\label{sec:Restless Legs Syndrome}
Restless legs syndrome (RLS, Willis-Ekbom disease) is a sensorimotor disorder characterized by an irresistible urge to move the legs, usually accompanied by uncomfortable sensations (creeping, crawling, pulling, aching, or tingling), that begins or worsens during periods of rest or inactivity, is partially or completely relieved by movement, occurs predominantly in the evening or night, and is not better explained by another condition. Prevalence is 5--10\% in adults of European descent, 2--4\% in Asian populations, and increases with age, affecting women 2:1 over men. The condition results from brain iron deficiency (reduced iron content in substantia nigra and striatal dopamine receptors) and dopaminergic dysfunction, with genetic associations (BTBD9, MEIS1, MAP2K5). Secondary RLS is caused by iron deficiency anemia, pregnancy, end-stage renal disease, peripheral neuropathy, and medications (SSRIs, antihistamines, antipsychotics). Clinical features include an irresistible urge to move the legs with uncomfortable sensations that worsen at rest and improve with movement. Diagnosis is clinical using the five essential IRLSSG diagnostic criteria; serum ferritin levels should be checked (goal >75 $\mu$g/L). Management includes non-pharmacologic measures (moderate exercise, avoidance of alcohol and caffeine in the evening), correction of iron deficiency, first-line pharmacotherapy with alpha-2-delta calcium channel ligands (gabapentin, pregabalin---preferred due to lower augmentation risk), and dopamine agonists (pramipexole, ropinirole, rotigotine) for refractory cases. Prognosis is chronic but manageable; augmentation is the most significant complication of dopaminergic therapy.

\subsection{Sleep-Related Bruxism}
\label{sec:Sleep-Related Bruxism}
Sleep-related bruxism is a sleep-related movement disorder characterized by repetitive jaw-muscle activity during sleep, manifesting as tooth grinding or clenching, often associated with rhythmic masticatory muscle activity (RMMA). Prevalence is 8--12\% in adults and 15--30\% in children, decreasing with age. The condition results from sleep-related activation of the brainstem masticatory central pattern generator, triggered by spontaneous microarousals and possibly modulated by autonomic activation; risk factors include anxiety, stress, sleep-disordered breathing, alcohol, smoking, and certain medications (SSRIs, amphetamines). Clinical features include audible tooth grinding, jaw pain or fatigue, temporal headaches, tooth wear, tooth sensitivity, and temporomandibular joint pain. Diagnosis is clinical; polysomnography with masseter EMG is the gold standard. Management includes occlusal splints for tooth protection, stress reduction, treatment of comorbid OSA, botulinum toxin injections for severe cases, and avoidance of exacerbating substances. Prognosis is generally manageable; severe untreated bruxism can result in tooth fracture and TMJ disorders.

\subsection{Sleep-Related Rhythmic Movement Disorder}
\label{sec:Sleep-Related Rhythmic Movement Disorder}
Sleep-related rhythmic movement disorder (RMD) is a sleep-related movement disorder involving stereotypical, repetitive movements of large muscle groups occurring primarily during the transition from wakefulness to sleep or during NREM sleep. The most common forms are head banging, head rolling, and body rocking. Prevalence is up to 60\% in infants (physiologic), declining to less than 1\% in adults. The condition results from unknown pathophysiology; in adults, it is associated with neurodevelopmental disorders, anxiety, and OSA. Clinical features include rhythmic movements at a frequency of 0.5--2 Hz lasting less than 15 minutes. Diagnosis is clinical; polysomnography with synchronized video recording is indicated when injury is suspected or to rule out seizure. Management is primarily reassurance and safety measures; clonazepam, hypnotics, or behavioral therapy may be used in persistent or injurious adult cases. Prognosis is benign; most children outgrow the condition.

